\documentclass[11pt]{article}

\usepackage[margin=1in]{geometry}
\usepackage{amsmath,amssymb,amsthm,mathtools}
\usepackage{microtype}
\usepackage{hyperref}
\usepackage{enumitem}

\hypersetup{
  colorlinks=true,
  linkcolor=blue,
  citecolor=blue,
  urlcolor=blue,
  pdftitle={Finite Forward-Iterate Schroeder Approximation from Geometric Richardson Extrapolation},
  pdfauthor={Wayne Baker}
}

\setlength{\parindent}{0pt}
\setlength{\parskip}{0.45em}
\setlength{\emergencystretch}{2em}

\newtheorem{theorem}{Theorem}[section]
\newtheorem{lemma}[theorem]{Lemma}
\newtheorem{proposition}[theorem]{Proposition}
\newtheorem{corollary}[theorem]{Corollary}
\theoremstyle{remark}
\newtheorem{remark}[theorem]{Remark}

\newcommand{\D}{\mathcal D}
\newcommand{\Lop}{\mathcal L}
\newcommand{\RR}{\mathbb R}
\newcommand{\qbinom}[3]{\begin{bmatrix}#1\\#2\end{bmatrix}_{#3}}

\title{\textbf{Finite Forward-Iterate Schr\"oder Approximation\
from Geometric Richardson Extrapolation}\\
\large Reciprocal Filters and Exact Paired Remainders}
\author{Wayne Baker\\\small Independent researcher}
\date{September 2026\\\small Revised theorem note v0.2.1}

\begin{document}
\maketitle

\begin{abstract}
Let \(\Phi\) be an odd analytic germ with \(\Phi(0)=0\) and
\(\Phi'(0)=1\), let \(H=\Phi^{-1}\), and fix \(m>1\).  The conjugated
multiplication map
\[
S_m(x)=\Phi\!\bigl(mH(x)\bigr)
\]
satisfies the classical Schr\"oder relation \(H\circ S_m=mH\).  For
\(y=H(x)\), consider the Richardson family
\[
T_y(h)=\frac{\Phi(hy)}{h}.
\]
At the geometric nodes \(h=m^q\), one has the exact identity
\[
T_y(m^q)=m^{-q}S_m^{\circ q}(x).
\]
Thus geometric Richardson extrapolation of \(T_y\) to \(h=0\) becomes,
in the original variable, a finite fixed-weight linear combination of
forward iterates of the repelling map \(S_m\).  This note develops that
realization and its reciprocal shrinking-scale form.  The resulting
composition and scale filters satisfy
\[
\ell_r(z)=z^r p_r(z^{-1}),
\qquad
p_r(z)=\prod_{j=1}^r\frac{m^{2j}-z}{m^{2j}-1},
\]
and give paired exact remainder identities.  In the generic case
\[
H(x)-P_{r+1}^{(m,\Phi)}(x)
=
(-1)^{r+1}a_{r+1}m^{r(r+1)}x^{2r+3}+O(x^{2r+5}),
\]
where \(\Phi(t)=t+\sum_{n\ge1}a_nt^{2n+1}\).  The Richardson filters,
geometric-node error factors, skipped-mode order jumps, and the underlying
Schr\"oder--Koenigs linearization are classical inputs; the emphasis here
is the explicit finite forward-iterate realization, its reciprocal pairing
with the shrinking-scale filter, and the exact remainder formulas in the
original variable.
\end{abstract}

\section{Purpose, claim boundary, and related work}

The construction in this note sits at the intersection of two classical
theories.  Schr\"oder--Koenigs linearization converts a nonlinear map near
a fixed point into multiplication in a conjugate coordinate, while
Richardson extrapolation removes prescribed asymptotic power modes by
finite weighted combinations.  Neither ingredient is claimed as new.

The observation developed here is deliberately narrower.  If
\(H\circ S_m=mH\) and \(y=H(x)\), then geometric Richardson data for
\[
T_y(h)=\frac{\Phi(hy)}{h}
\]
at the forward nodes \(h=1,m,\ldots,m^r\) are exactly evaluations of
finite forward iterates of \(S_m\) in the original variable.  Consequently
the Richardson extrapolant can be written without evaluating \(H\):
\[
P_{r+1}^{(m,\Phi)}(x)
=
\sum_{q=0}^r\lambda_{r,q}m^{-q}S_m^{\circ q}(x).
\]
The corresponding shrinking-scale Richardson filter is the reciprocal
polynomial \(\ell_r(z)=z^rp_r(z^{-1})\).  To the author's knowledge, this
finite forward-iterate realization and its reciprocal pairing with the
shrinking-scale filter have not previously been stated in this form.  This
is a literature-status statement, not an absolute priority claim.

Several parts of the calculation are standard and are used here as such.
Schr\"oder's equation, Koenigs linearization, and composition-operator
eigenfunctions are classical \cite{Schroeder1871,Koenigs1884}.  Targonski
studied Schr\"oder equations, iterated-function interpolation, and
Chebyshev generalizations well before the present construction
\cite{Targonski1957,Targonski1965,Targonski1967}.  Mira's 1982 contribution
on Schr\"oder equations and generalized Chebyshev polynomials is closely
adjacent to the map-construction side \cite{Mira1982}; the full text was
not available for direct inspection during this revision, so no priority
claim concerning its detailed formulas is made here.

On the extrapolation side, geometric-node Richardson and its error
analysis belong to the classical theory; see Joyce's survey,
Brezinski--Redivo-Zaglia, and Sidi; see especially Sidi, Chapter~1, for the Richardson framework and the classical Archimedean $\pi$ example \cite{Joyce1971,BrezinskiRedivo1991,Sidi2003}.
In particular, vanishing asymptotic coefficients and the resulting order
jumps are inherited Richardson phenomena rather than a separate novelty
claim.  The sine sequence \(n\sin(1/n)\) also appears as a standard
extrapolation example in the classical literature, so polygonal-sine or
polygonal-\(\pi\) acceleration is treated below only as historical context
and specialization.

Iserles developed constructive Schr\"oder approximation and a substantial
line of work on acceleration of functional iteration
\cite{Iserles1991,IserlesSaffLiu1992,Iserles1993,Iserles1994}.  In the 1993
paper the Schr\"oder function is constructed through Taylor-coefficient
recurrences and an auxiliary generating function.  His iteration argument uses
forward iterates of an attracting map, with $0<|\beta|<1$, either taken to the
Koenigs limit or, in the acceleration work, supplied to nonlinear
Hankel/Steffensen transforms.  The construction here instead uses finitely many
forward iterates of the repelling map $S_m$ with fixed linear weights.  The two
routes are therefore related but distinct.

The goal of this note is consequently modest and concrete: to formulate
that realization cleanly, derive exact paired remainders on the
composition and scale sides, and separate the general analytic-germ
statement from the special polynomial structure of the sine/Chebyshev
case.

Throughout the main theorem, \(m>1\) is an arbitrary real number.  Oddness
of \(\Phi\) restricts the spectral ladder to odd powers and the associated
Richardson expansion to even powers.  The underlying diagonalization does
not itself require oddness.  The restriction to odd integer \(m\) is
needed only in the sine specialization when one wants \(S_m\) to be a
polynomial.

\section{Odd analytic conjugacy and the forward-iterate realization}

Let
\begin{equation}
\Phi(t)
=
 t+\sum_{n\ge1}a_n t^{2n+1}
\label{eq:Phi}
\end{equation}
be an odd real-analytic germ at the origin, normalized by
\(\Phi'(0)=1\).  Let \(H=\Phi^{-1}\) denote its local analytic inverse.
For \(m>1\), define
\begin{equation}
S_m(x)=\Phi\!\bigl(mH(x)\bigr),
\qquad
(A_mf)(x)=\frac1m f(S_m(x)).
\label{eq:SmAm}
\end{equation}
Then
\begin{equation}
H(S_m(x))=mH(x).
\label{eq:Schroeder}
\end{equation}

\begin{lemma}[Diagonal action]
\label{lem:diag}
For every integer \(k\ge1\),
\begin{equation}
A_m(H^k)=m^{k-1}H^k.
\label{eq:fullDiag}
\end{equation}
In particular,
\begin{equation}
A_m(H^{2n+1})=m^{2n}H^{2n+1},
\qquad n\ge0.
\label{eq:oddDiag}
\end{equation}
\end{lemma}

\begin{proof}
Using \eqref{eq:Schroeder},
\[
A_m(H^k)
=
\frac1m H(S_m(x))^k
=
\frac1m(mH(x))^k
=
m^{k-1}H(x)^k.
\]
\end{proof}

\subsection{Richardson realization at forward geometric nodes}

Fix \(x\) sufficiently near the origin and write \(y=H(x)\).  Define
\begin{equation}
T_y(h):=\frac{\Phi(hy)}{h}.
\label{eq:Ty}
\end{equation}
By \eqref{eq:Phi},
\begin{equation}
T_y(h)
=
y+\sum_{n\ge1}a_n y^{2n+1}h^{2n}.
\label{eq:Tyexp}
\end{equation}
Thus the target value is \(T_y(0)=y=H(x)\), understood by analytic
continuation at \(h=0\).  At the geometric nodes \(h=m^q\),
\begin{equation}
\boxed{
T_y(m^q)
=
m^{-q}\Phi(m^qH(x))
=
m^{-q}S_m^{\circ q}(x)
=
A_m^q x.
}
\label{eq:Tyiterate}
\end{equation}
Choose a neighborhood \(U\) of the origin on which \(\Phi\) is analytic
and one-to-one, so that \(H\circ\Phi=\mathrm{id}\) on \(U\).  For a fixed
stage \(r\), the identity \eqref{eq:Tyiterate} is valid whenever
\(m^qH(x)\in U\) for every \(0\le q\le r\).  This condition ensures both
that all evaluations of \(\Phi\) are defined and that the inverse relation is
valid at every intermediate iterate.  Under this condition the finite
composition formula below is precisely geometric-node Richardson
extrapolation of \(T_y\) to \(h=0\), written back in the original variable.
The forward nodes grow away from zero; no limiting forward orbit is required.

\subsection{Reciprocal filters}

For \(r\ge0\), set
\begin{equation}
p_r(z)
=
\prod_{j=1}^r\frac{m^{2j}-z}{m^{2j}-1},
\qquad p_0(z)=1,
\label{eq:pr}
\end{equation}
and define the canonical finite iterate lift
\begin{equation}
P_{r+1}^{(m,\Phi)}(x)
:=
p_r(A_m)x.
\label{eq:Pdef}
\end{equation}

On the scale side define the reciprocal polynomial
\begin{equation}
\ell_r(z)
=
\prod_{j=1}^r\frac{m^{2j}z-1}{m^{2j}-1}.
\label{eq:ellr}
\end{equation}

\begin{lemma}[Reciprocal filter identity]
\label{lem:reciprocal}
For every \(r\ge0\),
\begin{equation}
\boxed{
\ell_r(z)=z^r p_r(z^{-1}).
}
\label{eq:reciprocal}
\end{equation}
\end{lemma}

\begin{proof}
Directly,
\[
z^r p_r(z^{-1})
=
\prod_{j=1}^r
\frac{z(m^{2j}-z^{-1})}{m^{2j}-1}
=
\prod_{j=1}^r
\frac{m^{2j}z-1}{m^{2j}-1}.
\]
\end{proof}

The identity \eqref{eq:reciprocal} is the algebraic source of the
composition/scale duality below.

\section{Exact composition-side remainder}

Because \(x=\Phi(H)\), equation \eqref{eq:Phi} becomes
\begin{equation}
x
=
H+\sum_{n\ge1}a_nH^{2n+1}.
\label{eq:xHexp}
\end{equation}
The filter \(p_r\) satisfies
\begin{equation}
p_r(1)=1,
\qquad
p_r(m^{2n})=0,
\quad 1\le n\le r.
\label{eq:roots}
\end{equation}

\begin{theorem}[Exact local active-mode remainder]
\label{thm:localExact}
Let \(R_\Phi\in(0,\infty]\) denote the radius of convergence of
\eqref{eq:Phi}.  For every \(r\ge0\), the canonical lift satisfies the
convergent germ identity below; its displayed series converges whenever
\(m^r|H(x)|<R_\Phi\).
\begin{equation}
\boxed{
P_{r+1}^{(m,\Phi)}
=
H+
\sum_{n\ge r+1}
a_n p_r(m^{2n})H^{2n+1}.
}
\label{eq:Pexact}
\end{equation}
Equivalently,
\begin{equation}
\boxed{
H-P_{r+1}^{(m,\Phi)}
=
-\sum_{n\ge r+1}
a_n p_r(m^{2n})H^{2n+1}.
}
\label{eq:localExact}
\end{equation}
\end{theorem}

\begin{proof}
Apply the finite operator polynomial \(p_r(A_m)\) to
\eqref{eq:xHexp}.  Lemma~\ref{lem:diag} gives
\[
p_r(A_m)x
=
p_r(1)H
+
\sum_{n\ge1}a_n p_r(m^{2n})H^{2n+1}.
\]
Now use \eqref{eq:roots}.
\end{proof}

Thus the usual leading defect is merely the first nonzero term in an
exact filtered tail.

\section{Finite iterate representation and exact intertwining}

Write
\begin{equation}
p_r(z)=\sum_{q=0}^r\lambda_{r,q}z^q.
\label{eq:lambdas}
\end{equation}
Since
\begin{equation}
A_m^q x
=
m^{-q}S_m^{\circ q}(x),
\label{eq:Aiterate}
\end{equation}
we have the finite iterate representation
\begin{equation}
\boxed{
P_{r+1}^{(m,\Phi)}(x)
=
\sum_{q=0}^r
\lambda_{r,q}m^{-q}S_m^{\circ q}(x).
}
\label{eq:finiteIterate}
\end{equation}

Fix a real target \(\Lambda\).  Choose the univalence neighborhood
\(U\) as above and take \(N>0\) sufficiently large that
\(m^q\Lambda/(m^rN)\in U\) for every \(0\le q\le r\); equivalently, all
intermediate points required by the orbit identity remain in the chosen
Schr\"oder coordinate chart.  Then define
\begin{equation}
G_N^{(\Phi,\Lambda)}
=
N\Phi\!\left(\frac{\Lambda}{N}\right).
\label{eq:Ggeneral}
\end{equation}
Let
\begin{equation}
(\D_mG)_N=G_{mN},
\qquad
\Lop_r^{(m)}
:=
\ell_r(\D_m)
=
\prod_{j=1}^r
\frac{m^{2j}\D_m-I}{m^{2j}-1}.
\label{eq:Ldef}
\end{equation}
When a genuinely discrete sequence is desired, one may restrict to
integer \(m\ge2\) and integer \(N\); no such arithmetic restriction is
needed for the analytic identity itself.

Set
\begin{equation}
x_K:=\Phi\!\left(\frac{\Lambda}{K}\right).
\label{eq:xK}
\end{equation}
Then
\begin{equation}
S_m(x_{mK})=x_K,
\qquad
S_m^{\circ q}(x_{m^rN})=x_{m^{r-q}N}.
\label{eq:orbit}
\end{equation}

\begin{theorem}[Exact Schr\"oder--Richardson intertwining]
\label{thm:intertwine}
For every \(r\ge0\) and every \(N\) satisfying the preceding
coordinate-domain condition,
\begin{equation}
\boxed{
\Lop_r^{(m)}G_N^{(\Phi,\Lambda)}
=
m^rN\,
P_{r+1}^{(m,\Phi)}
\!\left(
\Phi\!\left(\frac{\Lambda}{m^rN}\right)
\right).
}
\label{eq:intertwine}
\end{equation}
Consequently,
\begin{equation}
\boxed{
\Lambda-
\Lop_r^{(m)}G_N^{(\Phi,\Lambda)}
=
m^rN
\left[
H(x_{m^rN})-
P_{r+1}^{(m,\Phi)}(x_{m^rN})
\right].
}
\label{eq:defectIntertwine}
\end{equation}
\end{theorem}

\begin{proof}
Using \eqref{eq:finiteIterate} and \eqref{eq:orbit},
\begin{align*}
m^rN P_{r+1}^{(m,\Phi)}(x_{m^rN})
&=
\sum_{q=0}^r
\lambda_{r,q}
m^{r-q}N x_{m^{r-q}N}
\\
&=
\sum_{q=0}^r
\lambda_{r,q}
G_{m^{r-q}N}^{(\Phi,\Lambda)}.
\end{align*}
On the other hand, Lemma~\ref{lem:reciprocal} and
\eqref{eq:lambdas} give
\[
\ell_r(z)
=
z^rp_r(z^{-1})
=
\sum_{q=0}^r\lambda_{r,q}z^{r-q}.
\]
Therefore
\[
\Lop_r^{(m)}G_N^{(\Phi,\Lambda)}
=
\sum_{q=0}^r
\lambda_{r,q}
G_{m^{r-q}N}^{(\Phi,\Lambda)},
\]
which proves \eqref{eq:intertwine}.  Since
\[
H(x_{m^rN})=\frac{\Lambda}{m^rN},
\]
subtracting \eqref{eq:intertwine} from \(\Lambda\) gives
\eqref{eq:defectIntertwine}.
\end{proof}

The reciprocal identity \(\ell_r(z)=z^rp_r(z^{-1})\) therefore converts
the composition-side filter into the scale-side filter exactly.

\section{Exact sequence remainder}

The scale expansion corresponding to \eqref{eq:xHexp} is
\begin{equation}
G_N^{(\Phi,\Lambda)}
=
\Lambda+
\sum_{n\ge1}
a_n\Lambda^{2n+1}N^{-2n}.
\label{eq:Gexp}
\end{equation}

\begin{lemma}[Scale eigenvalues]
\label{lem:scaleEigen}
For every \(n\ge1\),
\begin{equation}
\boxed{
\Lop_r^{(m)}(N^{-2n})
=
m^{-2rn}p_r(m^{2n})N^{-2n}.
}
\label{eq:scaleEigen}
\end{equation}
\end{lemma}

\begin{proof}
Since \(\D_m(N^{-2n})=m^{-2n}N^{-2n}\), the multiplier contributed by
\(\Lop_r^{(m)}\) is
\[
\prod_{j=1}^r
\frac{m^{2j-2n}-1}{m^{2j}-1}.
\]
For each factor,
\[
m^{2j-2n}-1
=
m^{-2n}(m^{2j}-m^{2n}).
\]
Extracting \(m^{-2n}\) from all \(r\) factors gives
\eqref{eq:scaleEigen}.
\end{proof}

\begin{theorem}[Exact scale-side active-mode remainder]
\label{thm:sequenceExact}
For all sufficiently large \(N\),
\begin{equation}
\boxed{
\Lambda-
\Lop_r^{(m)}G_N^{(\Phi,\Lambda)}
=
-\sum_{n\ge r+1}
a_n p_r(m^{2n})m^{-2rn}
\Lambda^{2n+1}N^{-2n}.
}
\label{eq:sequenceExact}
\end{equation}
\end{theorem}

\begin{proof}
Apply \(\Lop_r^{(m)}\) termwise to the convergent expansion
\eqref{eq:Gexp}.  The constant mode is preserved, while
Lemma~\ref{lem:scaleEigen} gives the multiplier of each inverse-power
mode.  For \(1\le n\le r\), the factor \(p_r(m^{2n})\) vanishes.
\end{proof}

Equations \eqref{eq:localExact} and \eqref{eq:sequenceExact} display the
same active coefficients and the same spectral factor
\(p_r(m^{2n})\); only the reciprocal scale multiplier \(m^{-2rn}\)
separates the two sides.

\section{Skipped modes and exact first-surviving coefficients}

The order-jump mechanism in this section is the familiar Richardson
phenomenon that occurs when one or more nominal asymptotic coefficients
vanish.  The purpose here is not to claim a new superconvergence principle,
but to record its exact specialization to the present conjugacy and the
closed coefficient inherited from the geometric filter.

Assume the filtered tail is not identically zero, and define its first
active index
\begin{equation}
n_\star
:=
\min\{n\ge r+1:a_n\neq0\}.
\label{eq:nstar}
\end{equation}

For \(q>0\), define the Gaussian binomial coefficient by
\begin{equation}
\begin{bmatrix}L\\R\end{bmatrix}_q
:=
\prod_{j=1}^R
\frac{q^{L+1-j}-1}{q^j-1},
\qquad
0\le R\le L.
\label{eq:qbinom}
\end{equation}

\begin{lemma}[First-surviving filter value]
\label{lem:qbinom}
For every \(n_\star\ge r+1\),
\begin{equation}
\boxed{
p_r(m^{2n_\star})
=
(-1)^r m^{r(r+1)}
\begin{bmatrix}n_\star-1\\r\end{bmatrix}_{m^2}.
}
\label{eq:pqbinom}
\end{equation}
\end{lemma}

\begin{proof}
Starting from \eqref{eq:pr},
\begin{align*}
p_r(m^{2n_\star})
&=
\prod_{j=1}^r
\frac{m^{2j}-m^{2n_\star}}{m^{2j}-1}
\\
&=
(-1)^r m^{2\sum_{j=1}^rj}
\prod_{j=1}^r
\frac{m^{2(n_\star-j)}-1}{m^{2j}-1}.
\end{align*}
Since \(2\sum_{j=1}^rj=r(r+1)\), setting \(q=m^2\) turns the remaining
product into \(\bigl[\begin{smallmatrix}n_\star-1\\r\end{smallmatrix}\bigr]_q\).
\end{proof}

\begin{theorem}[Exact first-surviving coefficient]
\label{thm:super}
Let \(n_\star\) be given by \eqref{eq:nstar}.  Then
\begin{equation}
\boxed{
H(x)-P_{r+1}^{(m,\Phi)}(x)
=
(-1)^{r+1}a_{n_\star}
m^{r(r+1)}
\begin{bmatrix}n_\star-1\\r\end{bmatrix}_{m^2}
x^{2n_\star+1}
+
O(x^{2n_\star+3}).
}
\label{eq:localSuper}
\end{equation}
On the scale side,
\begin{equation}
\boxed{
\Lambda-
\Lop_r^{(m)}G_N^{(\Phi,\Lambda)}
=
(-1)^{r+1}a_{n_\star}
\begin{bmatrix}n_\star-1\\r\end{bmatrix}_{m^2}
m^{-r(2n_\star-r-1)}
\Lambda^{2n_\star+1}N^{-2n_\star}
+
O(N^{-2n_\star-2}).
}
\label{eq:sequenceSuper}
\end{equation}
\end{theorem}

\begin{proof}
In \eqref{eq:localExact}, the first surviving term is
\(-a_{n_\star} p_r(m^{2n_\star})H^{2n_\star+1}\).  Since \(H(x)=x+O(x^3)\),
\[
H^{2n_\star+1}=x^{2n_\star+1}+O(x^{2n_\star+3}).
\]
Substitution of \eqref{eq:pqbinom} proves \eqref{eq:localSuper}.

Likewise, the first surviving term of \eqref{eq:sequenceExact} is
\[
-a_{n_\star} p_r(m^{2n_\star})m^{-2rn_\star}
\Lambda^{2n_\star+1}N^{-2n_\star}.
\]
Using \eqref{eq:pqbinom} and
\[
r(r+1)-2rn_\star=-r(2n_\star-r-1)
\]
gives \eqref{eq:sequenceSuper}.
\end{proof}

\begin{corollary}[Generic paired defect]
\label{cor:generic}
If \(a_{r+1}\neq0\), then \(n_\star=r+1\) and
\begin{align}
H-P_{r+1}^{(m,\Phi)}
&=
(-1)^{r+1}a_{r+1}
m^{r(r+1)}x^{2r+3}
+O(x^{2r+5}),
\label{eq:genericLocal}
\\
\Lambda-
\Lop_r^{(m)}G_N^{(\Phi,\Lambda)}
&=
(-1)^{r+1}a_{r+1}
m^{-r(r+1)}
\Lambda^{2r+3}N^{-(2r+2)}
+O(N^{-(2r+4)}).
\label{eq:genericSequence}
\end{align}
\end{corollary}

Thus every missing coefficient
\(a_{r+1},a_{r+2},\ldots\) forces a corresponding automatic order jump.
In particular, \(a_{r+1}=0\) raises the scale-side order from
\(2r+2\) to at least \(2r+4\).

\section{Uniqueness and finite exact recovery}

The canonical filter \(p_r\) always exists, but jet agreement need not
characterize it uniquely when some low active modes are absent.

Let
\begin{equation}
\mathcal A_r
:=
\{n\in\{1,\ldots,r\}:a_n\neq0\},
\qquad
k_r:=|\mathcal A_r|.
\label{eq:activeSet}
\end{equation}

\begin{theorem}[Active-mode uniqueness criterion]
\label{thm:uniqueness}
Let \(q\) be a polynomial of degree at most \(r\), and set
\(Q=q(A_m)x\).  Then
\begin{equation}
Q-H=O(H^{2r+3})
\label{eq:jetTarget}
\end{equation}
if and only if
\begin{equation}
q(1)=1,
\qquad
q(m^{2n})=0
\quad(n\in\mathcal A_r).
\label{eq:activeConditions}
\end{equation}
The affine space of such polynomials has dimension
\begin{equation}
\boxed{r-k_r.}
\label{eq:dimension}
\end{equation}
Consequently, the jet-recovery filter is unique exactly when
\begin{equation}
\boxed{a_1a_2\cdots a_r\neq0,}
\label{eq:nondeg}
\end{equation}
in which case the unique solution is \(q=p_r\).
\end{theorem}

\begin{proof}
From \eqref{eq:xHexp} and Lemma~\ref{lem:diag},
\[
q(A_m)x
=
q(1)H+
\sum_{n\ge1}a_nq(m^{2n})H^{2n+1}.
\]
Therefore \eqref{eq:jetTarget} is equivalent to
\eqref{eq:activeConditions}.  The nodes
\[
1,
\qquad
m^{2n}\quad(n\in\mathcal A_r)
\]
are pairwise distinct.  Evaluation at these \(k_r+1\) nodes has full rank
on the \((r+1)\)-dimensional space of polynomials of degree at most \(r\).
Hence the solution affine space has dimension
\((r+1)-(k_r+1)=r-k_r\).  It is zero-dimensional exactly when
\(k_r=r\), equivalently \eqref{eq:nondeg}.  In that case the unique
polynomial has value one at \(1\) and zeros at
\(m^2,\ldots,m^{2r}\), hence is precisely \(p_r\).
\end{proof}

\begin{corollary}[Finite exact recovery]
\label{cor:exact}
If
\begin{equation}
a_n=0
\qquad(n\ge r+1),
\label{eq:finiteSupport}
\end{equation}
then
\begin{equation}
\boxed{
P_{r+1}^{(m,\Phi)}=H
}
\label{eq:exactLocal}
\end{equation}
as analytic germs, and
\begin{equation}
\boxed{
\Lop_r^{(m)}G_N^{(\Phi,\Lambda)}=\Lambda
}
\label{eq:exactSequence}
\end{equation}
for all sufficiently large \(N\).
\end{corollary}

\begin{proof}
Under \eqref{eq:finiteSupport}, the exact tails in
\eqref{eq:localExact} and \eqref{eq:sequenceExact} are empty.
\end{proof}

\section{Recursion and the polynomial-map corollary}

The canonical filters satisfy
\begin{equation}
p_j(z)
=
p_{j-1}(z)
\frac{m^{2j}-z}{m^{2j}-1}.
\label{eq:filterRecursion}
\end{equation}

\begin{proposition}[One-step canonical lift]
\label{prop:recursion}
For \(j\ge1\),
\begin{equation}
\boxed{
P_{j+1}^{(m,\Phi)}
=
\frac{m^{2j}P_j^{(m,\Phi)}-A_mP_j^{(m,\Phi)}}{m^{2j}-1}.
}
\label{eq:operatorRecursion}
\end{equation}
Equivalently,
\begin{equation}
\boxed{
P_{j+1}^{(m,\Phi)}
=
P_j^{(m,\Phi)}
+
\frac{mP_j^{(m,\Phi)}-P_j^{(m,\Phi)}\circ S_m}
{m^{2j+1}-m}.
}
\label{eq:functionRecursion}
\end{equation}
\end{proposition}

\begin{proof}
Equation \eqref{eq:operatorRecursion} is
\eqref{eq:filterRecursion} with \(z=A_m\), applied to \(x\).
Using
\[
A_mP_j^{(m,\Phi)}
=
\frac1m P_j^{(m,\Phi)}\circ S_m
\]
and rearranging yields \eqref{eq:functionRecursion}.
\end{proof}

For a general analytic germ, \(S_m\) and
\(P_{r+1}^{(m,\Phi)}\) need not be polynomials, so a universal degree
statement is not available.  It reappears exactly when the multiplication
map is polynomial.

\begin{corollary}[Polynomial multiplication maps]
\label{cor:degree}
Assume \(S_m\) is a polynomial of exact degree \(d\ge2\).  Then
\begin{equation}
\boxed{
\deg P_{r+1}^{(m,\Phi)}=d^r.
}
\label{eq:degree}
\end{equation}
\end{corollary}

\begin{proof}
Since \(\deg S_m^{\circ q}=d^q\), the deepest iterate in
\eqref{eq:finiteIterate} has degree \(d^r\), while all preceding terms
have smaller degree.  Its coefficient is
\[
\lambda_{r,r}
=
\frac{(-1)^r}{\prod_{j=1}^r(m^{2j}-1)}
\neq0,
\]
so the highest-degree term cannot cancel.
\end{proof}

For \(\Phi(t)=\sin t\) and odd integer \(m\), one has
\(S_m(x)=\sin(m\arcsin x)\), an odd polynomial of degree \(m\), and
\eqref{eq:degree} reduces to \(\deg P_{r+1}=m^r\).  More explicitly,
\[
S_m(x)=(-1)^{(m-1)/2}T_m(x),
\]
where \(T_m\) is the Chebyshev polynomial of the first kind.  Thus the
multiple-angle composition law and the exponential degree growth are
classical Chebyshev structure, not novelty claims.

\begin{remark}[Compositional economy versus polynomial degree]
In the sine case, \(P_{r+1}\) has degree \(m^r\), whereas the Taylor
polynomial of \(\arcsin x\) through degree \(2r+1\) already has error
\(O(x^{2r+3})\).  The forward-iterate construction is therefore not
polynomial-degree efficient.  Its structural advantage is different: the
stage-\(r\) approximant is assembled from \(r\) forward compositions of
the known map \(S_m\) and fixed scalar weights, without direct evaluation
of \(H\).  For a general analytic germ the admissible neighborhood also
depends on the largest forward node \(m^rH(x)\); no claim of fixed-\(x\)
convergence as \(r\to\infty\) is made.
\end{remark}

\section{Role of oddness and finite smoothness}

\begin{remark}[What oddness does]
The spectral identity \eqref{eq:fullDiag} does not require \(\Phi\) to be
odd.  If a general normalized local diffeomorphism has
\[
\Phi(H)=H+\sum_{k\ge2}c_kH^k,
\]
then \(H^k\) is still an \(A_m\)-eigenmode with eigenvalue
\(m^{k-1}\), while
\[
N\Phi(\Lambda/N)
=
\Lambda+
\sum_{k\ge2}c_k\Lambda^kN^{-(k-1)}.
\]
Oddness is therefore what restricts the active spectrum to
\(m^2,m^4,m^6,\ldots\) and the scale errors to
\(N^{-2},N^{-4},N^{-6},\ldots\).  It is essential to the clean even-power
Richardson ladder used in this note, but not to the underlying
Schr\"oder--Richardson mechanism itself.
\end{remark}

\begin{remark}[Regularity]
Analyticity is convenient because it makes
\eqref{eq:localExact} and \eqref{eq:sequenceExact} convergent identities.
At the purely algebraic level the entire construction is valid in the
ring of odd formal power series
\[
\Phi(t)\in t+t^3\RR[[t^2]].
\]
For a fixed generic stage \(r\), only finitely many Taylor coefficients
and sufficiently many derivatives are needed to obtain the leading
asymptotic formulas.  In particular, an odd \(C^{2r+5}\) local
diffeomorphism suffices for the generic defect through the stated remainder
order, by the standard Taylor formula with remainder applied to the finitely
many compositions occurring at that stage.
\end{remark}

\section{Unified theorem}

The preceding results may be compressed as follows.

\begin{theorem}[Finite forward-iterate Richardson principle]
\label{thm:master}
Let \(\Phi\) be an odd analytic germ satisfying
\[
\Phi(0)=0,
\qquad
\Phi'(0)=1,
\qquad
\Phi(t)=t+\sum_{n\ge1}a_nt^{2n+1},
\]
let \(H=\Phi^{-1}\), and let \(m>1\).  Define \(S_m,A_m,p_r,P_{r+1}\),
\(\ell_r\), \(G_N^{(\Phi,\Lambda)}\), and \(\Lop_r^{(m)}\) by
\eqref{eq:SmAm}, \eqref{eq:pr}, \eqref{eq:Pdef}, \eqref{eq:ellr},
\eqref{eq:Ggeneral}, and \eqref{eq:Ldef}.  Then:
\begin{enumerate}[label=(\alph*)]
\item \(A_m(H^{2n+1})=m^{2n}H^{2n+1}\), and
      \(\ell_r(z)=z^rp_r(z^{-1})\).
\item The exact local remainder is
\[
H-P_{r+1}^{(m,\Phi)}
=
-\sum_{n\ge r+1}
a_np_r(m^{2n})H^{2n+1}.
\]
\item The exact finite-stage intertwining is
\[
\Lop_r^{(m)}G_N^{(\Phi,\Lambda)}
=
m^rN P_{r+1}^{(m,\Phi)}
\!\left(\Phi\!\left(\frac{\Lambda}{m^rN}\right)\right).
\]
\item The exact scale-side remainder is
\[
\Lambda-\Lop_r^{(m)}G_N^{(\Phi,\Lambda)}
=
-\sum_{n\ge r+1}
a_np_r(m^{2n})m^{-2rn}
\Lambda^{2n+1}N^{-2n}.
\]
\item If \(n_\star=\min\{n\ge r+1:a_n\neq0\}\), then the first surviving
      local and scale coefficients are given by
      \eqref{eq:localSuper} and \eqref{eq:sequenceSuper}.  In particular,
      missing active coefficients produce automatic superconvergence.
\item Jet recovery through order \(2r+1\) uniquely determines the
      canonical degree-\(r\) filter if and only if
      \(a_1\cdots a_r\neq0\).  In general the family of admissible filters
      has affine dimension \(r-k_r\).
\item If \(a_n=0\) for all \(n\ge r+1\), then both the local and sequence
      recovery are exact after finitely many stages.
\end{enumerate}
\end{theorem}

\section{Three immediate specializations}

\subsection{Sine/arcsine family}

Take \(\Phi(t)=\sin t\).  Then \(H(x)=\arcsin x\) and
\[
a_n=\frac{(-1)^n}{(2n+1)!}.
\]
Every active coefficient is nonzero.  Hence the canonical filter is unique
at every stage, and Corollary~\ref{cor:generic} gives
\[
\arcsin x-P_{r+1}^{(m,\sin)}(x)
=
\frac{m^{r(r+1)}}{(2r+3)!}x^{2r+3}+O(x^{2r+5}),
\]
while
\[
\Lambda-\Lop_r^{(m)}
\left[N\sin\!\left(\frac{\Lambda}{N}\right)\right]
=
\frac{\Lambda^{2r+3}}
{(2r+3)!\,m^{r(r+1)}}N^{-(2r+2)}
+O(N^{-(2r+4)}).
\]
If \(m\) is an odd integer, then
\(S_m(x)=\sin(m\arcsin x)\) is a polynomial of degree \(m\), so
Corollary~\ref{cor:degree} gives
\[
\deg P_{r+1}^{(m,\sin)}=m^r.
\]
For such odd integer \(m\), the multiple-angle composition
law gives
\[
S_m^{\circ q}(x)=\sin\!\bigl(m^q\arcsin x\bigr),
\qquad x\in[-1,1].
\]
Hence the finite-stage remainder identity extends from a germ identity to
every \(x\in[-1,1]\):
\[
\boxed{
\arcsin x-P_{r+1}^{(m,\sin)}(x)
=
-\sum_{n\ge r+1}
\frac{(-1)^n}{(2n+1)!}
 p_r(m^{2n})\,(\arcsin x)^{2n+1}.
}
\]
This global finite-stage identity should not be confused with convergence
of \(P_{r+1}(x)\) for fixed nonzero \(x\) as \(r\to\infty\), which is not
asserted.

\subsection{Finite exact recovery after one stage}

Let
\[
\Phi(t)=t+ct^3.
\]
Then the only nontrivial active coefficient is \(a_1=c\).  For \(r=1\),
Theorem~\ref{thm:localExact} has an empty tail, so
\[
P_2^{(m,\Phi)}=H
\]
exactly as germs.  Likewise
\[
G_N^{(\Phi,\Lambda)}
=
\Lambda+c\Lambda^3N^{-2},
\]
and the one-stage scale filter gives
\[
\frac{m^2G_{mN}^{(\Phi,\Lambda)}-G_N^{(\Phi,\Lambda)}}{m^2-1}
=
\Lambda
\]
exactly.

\subsection{A one-stage superconvergent gap}

Suppose
\[
\Phi(t)
=
t+a_1t^3+a_3t^7+O(t^9),
\qquad a_3\neq0,
\]
so the coefficient \(a_2\) is absent.  With \(r=1\), the first unfiltered
active index is \(n_\star=3\).  Since
\[
\begin{bmatrix}2\\1\end{bmatrix}_{m^2}=1+m^2,
\]
Theorem~\ref{thm:super} gives
\[
H-P_2^{(m,\Phi)}
=
a_3m^2(1+m^2)x^7+O(x^9),
\]
and
\[
\Lambda-\Lop_1^{(m)}G_N^{(\Phi,\Lambda)}
=
a_3(1+m^2)m^{-4}\Lambda^7N^{-6}+O(N^{-8}).
\]
Thus the nominal one-stage scale order \(N^{-4}\) jumps automatically to
\(N^{-6}\) because the first unfiltered mode is missing.

\section{Conclusion}

The central identity of this note is elementary once the appropriate
Richardson family is exposed:
\[
T_y(m^q)
=
m^{-q}S_m^{\circ q}(x),
\qquad y=H(x).
\]
It identifies a geometric Richardson extrapolant in the Schr\"oder
coordinate with a fixed-weight linear combination of finite forward
iterates of the original repelling map.  The reciprocal relation
\[
\ell_r(z)=z^rp_r(z^{-1})
\]
then connects this forward realization to the corresponding shrinking-scale
filter.

The extrapolation mechanism, geometric-node weights, skipped-mode order
jumps, and Schr\"oder--Koenigs diagonalization are classical.  The point of
the present formulation is instead to make the finite forward-iterate
realization explicit and to record the exact paired tails
\[
H-P_{r+1}^{(m,\Phi)}
=
-\sum_{n\ge r+1}a_np_r(m^{2n})H^{2n+1}
\]
and
\[
\Lambda-\Lop_r^{(m)}G_N^{(\Phi,\Lambda)}
=
-\sum_{n\ge r+1}
a_np_r(m^{2n})m^{-2rn}\Lambda^{2n+1}N^{-2n}.
\]
In the generic case the first formula gives
\[
H-P_{r+1}^{(m,\Phi)}
=
(-1)^{r+1}a_{r+1}m^{r(r+1)}x^{2r+3}
+O(x^{2r+5}),
\]
while the reciprocal scale side carries the corresponding
\(m^{-r(r+1)}\) factor.

For the sine family with odd integer \(m\), the forward maps are
Chebyshev multiple-angle polynomials and the finite-stage remainder is global
on \([-1,1]\), but
the construction is not degree-efficient relative to Taylor approximation.
Its value is compositional: the approximant uses only evaluations of the
known forward map and fixed weights.  This distinction, together with the
paired exact remainders, is the appropriate scope of the present note.

\begin{thebibliography}{99}
\raggedright

\bibitem{Schroeder1871}
E.~Schr\"oder,
\newblock \emph{\"Uber iterirte Functionen},
\newblock Mathematische Annalen \textbf{3} (1871), 296--322.

\bibitem{Koenigs1884}
G.~Koenigs,
\newblock \emph{Recherches sur les int\'egrales de certaines \'equations fonctionnelles},
\newblock Annales scientifiques de l'\'Ecole Normale Sup\'erieure, 3e s\'erie,
\textbf{1} (1884), 3--41.

\bibitem{Targonski1957}
G.~Targonski,
\newblock \emph{Interpolation durch Reihen iterierter Funktionen},
\newblock Commentationes Physico-Mathematicae \textbf{20}, no.~9 (1957), 3--12.

\bibitem{Targonski1965}
G.~I.~Targonski,
\newblock \emph{Schr\"oder's Equation and a Generalization of the Chebyshev Polynomials},
\newblock Transactions of the New York Academy of Sciences \textbf{27} (1965), 600--606.

\bibitem{Targonski1967}
G.~I.~Targonski,
\newblock \emph{Seminar on Functional Operators and Equations},
\newblock Lecture Notes in Mathematics, vol.~33, Springer, 1967.

\bibitem{Mira1982}
C.~Mira,
\newblock \emph{\'Equation de Schr\"oder et solutions des r\'ecurrences.
G\'en\'eralisation des polyn\^omes de Tchebycheff},
\newblock in \emph{Th\'eorie de l'It\'eration et ses Applications},
Colloques Internationaux du CNRS, no.~332, Toulouse (1982), 35--43.

\bibitem{Joyce1971}
D.~C.~Joyce,
\newblock \emph{Survey of Extrapolation Processes in Numerical Analysis},
\newblock SIAM Review \textbf{13} (1971), 435--490.
\newblock doi:10.1137/1013092.

\bibitem{BrezinskiRedivo1991}
C.~Brezinski and M.~Redivo Zaglia,
\newblock \emph{Extrapolation Methods: Theory and Practice},
\newblock Studies in Computational Mathematics, vol.~2, North-Holland/Elsevier, 1991.

\bibitem{Sidi2003}
A.~Sidi,
\newblock \emph{Practical Extrapolation Methods: Theory and Applications},
\newblock Cambridge Monographs on Applied and Computational Mathematics, vol.~10,
Cambridge University Press, 2003.
\newblock doi:10.1017/CBO9780511546815.

\bibitem{Iserles1991}
A.~Iserles,
\newblock \emph{Complex Dynamics of Convergence Acceleration},
\newblock IMA Journal of Numerical Analysis \textbf{11} (1991), 205--240.
\newblock doi:10.1093/imanum/11.2.205.

\bibitem{IserlesSaffLiu1992}
A.~Iserles, E.~B.~Saff, and X.~Liu,
\newblock \emph{On Zeros of Hankel Determinants with Iterated Polynomial Entries},
\newblock IMA Journal of Numerical Analysis \textbf{12} (1992), 387--403.
\newblock doi:10.1093/imanum/12.3.387.

\bibitem{Iserles1993}
A.~Iserles,
\newblock \emph{A Constructive Approach to the Schr\"oder Equation},
\newblock Journal of Computational and Applied Mathematics \textbf{46} (1993), 301--314.
\newblock doi:10.1016/0377-0427(93)90304-T.

\bibitem{Iserles1994}
A.~Iserles,
\newblock \emph{Convergence Acceleration as a Dynamical System},
\newblock Applied Numerical Mathematics \textbf{15} (1994), 101--121.
\newblock doi:10.1016/0168-9274(94)00020-4.

\end{thebibliography}

\end{document}
